\section{Delta Kinematics}
\label{sec:kinematics}

Each cell is a translational delta driven by three prismatic actuators whose
axes are vertical and mutually parallel. The parallelogram limbs hold the
moving platform parallel to the base, so the platform translates in three
degrees of freedom without rotating and each limb reduces to a single distance
constraint.

Let the base frame be centred on the base-triangle centroid with
$\hat{\mathbf{z}}$ along the actuator axes, and index the limbs by azimuth
$\theta_i\in\{\pi/2,\,-\pi/6,\,7\pi/6\}$, $\hat{\mathbf{u}}_i=[\cos\theta_i,\
\sin\theta_i]^{\!\top}$. Writing $R_b$ and $R_p$ for the base and platform
circumradii, $q_i$ for the prismatic coordinates and $\mathbf{x}=[x,y,z]^{\!\top}$
for the platform centre, the carriage and platform joints of limb $i$ are
$\mathbf{c}_i=R_b[\hat{\mathbf{u}}_i^{\!\top},0]^{\!\top}+q_i\hat{\mathbf{z}}$ and
$\mathbf{p}_i=\mathbf{x}+R_p[\hat{\mathbf{u}}_i^{\!\top},0]^{\!\top}$, and for
limbs of length $L$ the loop closure is
\begin{equation}
\label{eq:closure}
\|\mathbf{p}_i-\mathbf{c}_i\|=L,\qquad i=1,2,3 .
\end{equation}

Base anchor and platform joint share the azimuth $\theta_i$, so their
difference collapses onto $\hat{\mathbf{u}}_i$ and \eqref{eq:closure} separates
into three scalar equations $\rho_i^2+(z-q_i)^2=L^2$, with
$\rho_i=\|[x,y]^{\!\top}+(R_p-R_b)\hat{\mathbf{u}}_i\|$. The inverse kinematics
is therefore decoupled and closed-form,
\begin{equation}
\label{eq:ik}
q_i=z-\sqrt{L^2-\rho_i^2},
\end{equation}
the negative branch placing each carriage below its platform joint, which is
the assembly mode of the prototype. A pose is reachable when $\rho_i\le L$ and
$q_i\in[q_{\min},q_{\max}]$ for every limb.

The forward map follows by absorbing the platform offset into the carriage
position, $\mathbf{s}_i=\mathbf{c}_i-R_p[\hat{\mathbf{u}}_i^{\!\top},0]^{\!\top}$,
which turns \eqref{eq:closure} into $\|\mathbf{x}-\mathbf{s}_i\|=L$: the
platform centre lies at the intersection of three spheres of equal radius. This
trilateration admits the standard closed-form solution, and of its two mirrored
roots the one placing the platform above the plane of the $\mathbf{s}_i$ is
taken, consistent with \eqref{eq:ik}.

The prototype uses $s_b=21.39$\,mm, $s_p=12.00$\,mm and $L=56.00$\,mm, with
$93$\,mm of prismatic travel; the tool tip sits $24$\,mm above the platform
plane along its (invariant) normal, and is reported in place of $\mathbf{x}$ by
both maps. The resulting workspace spans $93$\,mm vertically and reaches a
maximum radius of $53$\,mm.
